%=====================================================================
% FRONT MATTER
%=====================================================================
\pagenumbering{roman}
\structuralfigures{F.}
\thispagestyle{empty}

\begin{center}
\vspace*{1.0cm}

{\color{secondary}\rule{0.85\textwidth}{2pt}}\\[6pt]
{\large\textsc{\color{secondary}Bachelor of Science in Information Technology}}\\[18pt]

{\Huge\bfseries\color{primary}ANALYSIS OF ALGORITHMS}\\[10pt]
{\LARGE\bfseries\color{primary}Correctness, Cost, and the Limits of Efficiency}\\[16pt]

{\color{secondary}\rule{0.85\textwidth}{2pt}}\\[16pt]

{\Large\textsc{Unified Lecture Handout}}\\[6pt]
{\large\color{gray!70!black}Complete Course in One Volume $\bullet$ Fall 2026}\\[24pt]

% --- the claim the whole book rests on ---
\begin{tikzpicture}
  \node[rounded corners=4pt, draw=primary, line width=1.1pt, fill=primary!6,
        text=primary, align=center, font=\small\bfseries,
        minimum width=11.6cm, minimum height=1.5cm]
       {\faIcon{stopwatch}\\[2pt]Two questions about every algorithm:
        is it right, and what does it cost?};
\end{tikzpicture}\\[10pt]
{\small\itshape\color{gray!70!black}And a third the first two do not answer:
is there a better one, and how would you know?}

\vspace{1.0cm}

\begin{tcolorbox}[enhanced, width=0.86\textwidth, colback=lightbg,
  colframe=primary, boxrule=0pt, leftrule=4pt, arc=2pt, top=6pt, bottom=6pt]
\small
\textbf{\color{primary}What this volume covers.} Every clause of the HEC 2023
outline for \emph{Analysis of Algorithms}, in the order the outline gives them,
with the published learning outcomes traced chapter by chapter in the Course
Specification overleaf.

\smallskip
Under HEC 2023 this is a \textbf{computing core} course --- compulsory in all
ten BS computing degrees, this one included --- taken in Semester 7, worth
3 (3--0) credit hours, with \emph{Data Structures} as its prerequisite. The
2023 outline is word for word the 2017 outline, under which the course was
BSCS-only and absent from BSIT.

\smallskip
About a third of that outline you have already met in Data Structures.
Appendix~F draws the line: that course owns the structures, this one owns the
\emph{analysis} and the design paradigms.
\end{tcolorbox}

\vspace{0.9cm}
{\small\color{gray!60!black}Six Parts $\bullet$ Twenty-Three Chapters $\bullet$ Seven Appendices}
\vspace{1cm}
\end{center}

\newpage

%=====================================================================
\chapter*{How to Use This Handout}
\addcontentsline{toc}{chapter}{How to Use This Handout}
\markboth{How to Use This Handout}{}

Most of what this course teaches is not an algorithm. It is a way of arguing
about one.

By Semester 7 you can already write a program that sorts. What you cannot yet
do reliably is say \emph{why} it is correct for every input rather than the
ones you tried, \emph{what} it will cost when the input is a thousand times
larger, and \emph{whether} anything could do better. Those three questions ---
correctness, cost, and the limit --- are the whole subject, and each has a
technique behind it: the loop invariant, the asymptotic bound, and the lower
bound.

\section*{The six parts}

\textbf{Part~I --- Analysing an Algorithm} (Chapters 1--4) builds the
vocabulary. What ``efficient'' means and against what model; best, expected and
worst case; the five asymptotic notations defined properly, not just $O$; and
the loop invariant, which is how a correctness claim is proved rather than
asserted.

\textbf{Part~II --- Recursion and Divide-and-Conquer} (Chapters 5--9) is the
machinery. Recurrences solved three ways, the master theorem and where it does
not apply, brute force taken seriously as a baseline, then merge sort and
quicksort analysed in full.

\textbf{Part~III --- Sorting, Selection and the Structures Beneath}
(Chapters 10--13) establishes the first lower bound in the book --- no
comparison sort can beat $n \lg n$ --- and then shows what buying extra
assumptions gets you. Chapters 12 and 13 revisit heaps, search trees and
hashing from Data Structures, for their \emph{analysis} rather than their
implementation.

\textbf{Part~IV --- The Design Paradigms} (Chapters 14--17) is the part you will
use most after the examination: greedy, dynamic programming, and the amortised
analysis without which several earlier claims cannot be stated honestly.

\textbf{Part~V --- Graph Algorithms} (Chapters 18--20) applies all of it, and is
where sparsity starts changing every bound you write.

\textbf{Part~VI --- Strings, Limits and Practice} (Chapters 21--23) closes with
string matching, the complexity classes, and your project.

\begin{conceptbox}
Why the lower bound in Chapter 10 matters more than the algorithms around it.
Everything before it answers ``how fast is this?''. Chapter 10 answers ``how
fast could \emph{anything} be?'', and that is a different kind of statement ---
it is about the problem, not about any program. Once you have seen one proved,
Chapter 22's claim that we do not know whether $\mathrm{P} = \mathrm{NP}$ stops
being a slogan and becomes a recognisable open question of the same shape.
\end{conceptbox}

\section*{The furniture of every chapter}

\rowcolors{2}{white}{lightbg}
\begin{longtable}{@{}L{4.6cm}L{10.6cm}@{}}
\toprule
\rowcolor{primary!15}
\textbf{Box} & \textbf{What it is for} \\
\midrule
\endhead
\faIcon{bullseye}~\textbf{Outcomes} & What you will be able to \emph{do}. Bloom
verbs, not topics. \\
\faIcon{link}~\textbf{Before you start} & Which earlier chapters this one stands
on. \\
\faIcon{tags}~\textbf{Key terms} & The vocabulary introduced; all of it is in
the glossary of Appendix~D. \\
\faIcon{book}~\textbf{Definition} & The canonical statement. In this subject the
definitions are load-bearing: most wrong answers in the examination are a
misremembered definition, not a failed calculation. \\
\faIcon{award}~\textbf{Theorem}, \faIcon{puzzle-piece}~\textbf{Lemma},
\faIcon{angle-double-right}~\textbf{Corollary} & A result, numbered within the
chapter so the text can say ``by Lemma 6.2''. \\
\emph{Proof.} & Deliberately not in a box --- it is prose and should be read as
prose. Every proof here is one you could reproduce; the ones that are not are
cited instead. \\
\faIcon{list-ol}~\textbf{Algorithm} & Pseudocode, numbered by line, so the
analysis can point at line 4. \\
\faIcon{lightbulb}~\textbf{Concept} & The judgement call --- what an experienced
practitioner would decide, and why. \\
\faIcon{exclamation-triangle}~\textbf{Warning} & Something that will bite you,
usually in the examination. \\
\faIcon{calculator}~\textbf{Worked example} & Real arithmetic, every step shown.
Do them with a pen. \\
\faIcon{exclamation-circle}~\textbf{Common pitfalls} & The mistakes that are
actually made. \\
\faIcon{flask}~\textbf{Try it yourself} & The laboratory: C++ you compile, run
and \emph{time}. \\
\faIcon{question-circle}~\textbf{Checkpoint} & Three questions. Attempt them
before reading Appendix~C. \\
\faIcon{clipboard-check}~\textbf{Chapter summary} & What to retain. \\
\faIcon{pen-fancy}~\textbf{Review questions} & Examination practice: multiple
choice, short answer, applied. \\
\bottomrule
\end{longtable}

\begin{alertbox}[The measured strand, and why it is here]
Every chapter also carries at least one \emph{measurement}: the algorithm
compiled, run, and timed, with the figures printed in the book.

\smallskip
This is not decoration. Asymptotic analysis deliberately discards constant
factors, and a student who only ever sees the asymptotics draws two wrong
conclusions --- that an $\bigO{n\lg n}$ algorithm always beats an $\bigO{n^2}$
one, and that the analysis is therefore unreliable. Both are cured by seeing
where the crossover actually falls. Insertion sort really does beat merge sort
on small inputs, and the asymptotics really do win by the time it matters.

\smallskip
Every figure in those boxes was produced by the file named beside it, in
\texttt{code/}, with the compiler and flags recorded in Appendix~A. Rerun them.
Your numbers will differ from the book's; the \emph{shape} should not.
\end{alertbox}

\newpage

%=====================================================================
\chapter*{Course Specification}
\addcontentsline{toc}{chapter}{Course Specification}
\markboth{Course Specification}{}

This page is the course file. It states what HEC requires of a course called
\emph{Analysis of Algorithms} and shows where each requirement is met. Three
editions are in force in different programmes, so all three are traced.

\section*{Course particulars}

\rowcolors{2}{white}{lightbg}
\begin{longtable}{@{}L{4.2cm}L{11.0cm}@{}}
\toprule
\rowcolor{primary!15}
\textbf{Item} & \textbf{Value} \\
\midrule
\endhead
Course title & Analysis of Algorithms (HEC 2023). HEC 2017 and HEC 2025 both
call it \emph{Design and Analysis of Algorithms} \\
Programme & BS Information Technology. Under HEC 2023 this is a
\textbf{computing core} course --- compulsory in all ten BS computing degrees
--- in Semester~7. Under HEC 2017 it was a BSCS core course and
\emph{absent from BSIT entirely}; under HEC 2025 it is a BSCS major core in
Semester~IV \\
Credit hours & 3 (3--0): three credit hours of lectures, no separate laboratory
credit. The laboratory work in this handout is therefore done in lecture time
and as homework \\
Prerequisite & \emph{Data Structures} (HEC 2023). HEC 2017 required
\emph{Data Structures and Algorithms}; HEC 2025 publishes no prerequisites at
all \\
Duration & Sixteen teaching weeks, midterm in Week~8, project due Week~16
(Appendix~G) \\
Course introduction & The HEC 2023 statement, verbatim: ``Detailed study of the
basic notions of the design of algorithms and the underlying data structures.
Several measures of complexity are introduced. Emphasis on the structure,
complexity, and efficiency of algorithms.'' \\
Set texts & As published: Cormen, Leiserson, Rivest and Stein,
\emph{Introduction to Algorithms}; Kleinberg and Tardos, \emph{Algorithm
Design}; Sedgewick and Wayne, \emph{Algorithms}. See the note below on the
edition, and Appendix~E \\
\bottomrule
\end{longtable}

\begin{alertbox}[Four defects in the published documents]
Recorded here so that a reader comparing this handout with the HEC booklet is
not confused, and so that the department's own course file does not inherit
them.

\smallskip
\textbf{1. The 2023 CLO labels are mis-aligned with their own text.} In the
2023 booklet \texttt{CLO-1} sits against the first half of outcome~1 and
\texttt{CLO-2} against its continuation, leaving 3, 5 and 7 dangling. The 2017
listing of the same eight outcomes is correctly numbered, and \emph{that} is
the numbering used below.

\smallskip
\textbf{2. The Bloom's Taxonomy column is empty in both editions}, for all
eight outcomes --- although the neighbouring Data Structures table fills it in.
The levels in the table below are therefore \emph{assigned by this department},
and are marked as such.

\smallskip
\textbf{3. ``Corman''} is the booklets' misspelling of \textbf{Cormen}, in both
2017 and 2023.

\smallskip
\textbf{4. The 3rd edition (2009) is cited; the 4th edition (2022) is
current.} Appendix~E cites the 4th and keeps the 3rd beside it, because
chapter numbers moved between them.
\end{alertbox}

\section*{Coverage of the HEC 2023 course outline}

The published outline is a single run-on paragraph. Each of its clauses is
traced below, in the order HEC writes them. Nothing is omitted.

\rowcolors{2}{white}{lightbg}
\begin{longtable}{@{}L{6.4cm}L{2.0cm}L{6.6cm}@{}}
\toprule
\rowcolor{primary!15}
\textbf{HEC 2023 clause} & \textbf{Chapter} & \textbf{Where exactly} \\
\midrule
\endhead
``Introduction; role of algorithms in computing'' & 1 & What an algorithm is,
why the question is asymptotic, and the model of computation the whole book
assumes \\
``Analysis on nature of input and size of input'' & 2 & What ``size'' means for
an array, a graph, a matrix and a number; best, expected and worst case \\
``Asymptotic notations; Big-O, Big $\Omega$, Big $\Theta$, little-o,
little-$\omega$'' & 3 & All five defined formally and proved from the
definitions \\
``Sorting Algorithm analysis'' & 8, 9, 10 & Merge sort, quicksort, and the
lower bound that bounds them both \\
``loop invariants'' & 4 & Initialisation, maintenance, termination; insertion
sort proved correct \\
``Recursion and recurrence relations'' & 5, 6 & Substitution, recursion trees,
and the master theorem \\
``Algorithm Design Techniques'' & 7--9, 14--16 & The four the outline names,
each with the property that makes it applicable \\
``Brute Force Approach'' & 7 & Taken seriously: the baseline every other
technique must beat, and the cases where it wins \\
``Divide-and-conquer approach; Merge, Quick Sort'' & 8, 9 & Both, in full, with
the average-case analysis of quicksort \\
``Greedy approach'' & 14 & Greedy-choice and optimal substructure, proved by
exchange argument \\
``Dynamic programming; Elements of Dynamic Programming'' & 15, 16 & The two
elements as their own chapter, then the classic problems \\
``Search trees'' & 12 & Revisited for analysis: why the height is what it is,
and what that costs \\
``Heaps'' & 12 & Likewise, with the priority-queue bounds Chapters 19 and 20
depend on \\
``Hashing'' & 13 & Expected versus worst case, load factor, and what
``constant time'' actually promises \\
``Graph algorithms, shortest paths, sparse graphs'' & 18, 20 & Representation
and sparsity first, because sparsity changes every bound in the part; then
single-source and all-pairs \\
``String matching'' & 21 & Naive, Rabin--Karp and Knuth--Morris--Pratt \\
``Introduction to complexity classes'' & 22 & P, NP, NP-hard and NP-complete,
with one reduction worked in full \\
\bottomrule
\end{longtable}

\section*{Course learning outcomes}

The eight outcomes HEC publishes, in the 2017 numbering (see defect 1 above).
Bloom levels are assigned by this department, because the published tables
leave the column blank.

\rowcolors{2}{white}{lightbg}
\begin{longtable}{@{}L{1.3cm}L{7.6cm}L{2.1cm}L{3.8cm}@{}}
\toprule
\rowcolor{primary!15}
\textbf{CLO} & \textbf{By the end of the course the student will be able to} &
\textbf{Bloom} & \textbf{Chapters} \\
\midrule
\endhead
CLO-1 & Explain what is meant by ``best'', ``expected'' and ``worst'' case
behaviour of an algorithm & C2 Understand & 2, 9 \\
CLO-2 & Identify the characteristics of data, or other conditions or
assumptions, that lead to different behaviours & C4 Analyse & 2, 9, 13, 18 \\
CLO-3 & Determine informally the time and space complexity of simple
algorithms & C3 Apply & 1, 3, 17, and every chapter's worked example \\
CLO-4 & List and contrast standard complexity classes & C2 Understand & 22 \\
CLO-5 & Use big~$O$, $\Omega$ and $\Theta$ notation formally to give asymptotic
bounds on time and space complexity & C3 Apply & 3, 5, 6 \\
CLO-6 & Use the strategies --- brute force, greedy, divide-and-conquer and
dynamic programming --- to solve an appropriate problem & C3 Apply &
7, 8, 9, 14, 15, 16 \\
CLO-7 & Solve problems using graph algorithms, including single-source and
all-pairs shortest paths and at least one minimum spanning tree algorithm &
C3 Apply & 18, 19, 20 \\
CLO-8 & Trace and/or implement a string-matching algorithm & C3 Apply & 21 \\
\bottomrule
\end{longtable}

\begin{conceptbox}[One chapter the outline does not name, and why it is here]
\textbf{Chapter 19, Minimum Spanning Trees}, appears nowhere in the outline
paragraph --- and \textbf{CLO-7 requires it by name}: ``at least one minimum
spanning tree algorithm''. The outline and the outcomes disagree, and this
handout follows the outcome, because that is what a student is assessed
against.

\smallskip
Two further chapters go beyond the letter of the outline and are justified
rather than smuggled in. \textbf{Chapter 11 (Selection)} is the natural
companion to Chapter 10's lower bound and costs one lecture. \textbf{Chapter 17
(Amortised Analysis)} is needed before Chapter 13's hash-table claim or
Chapter 19's union-find can be stated honestly; without it, ``constant time''
in those chapters would have to be asserted rather than proved, and CLO-3 asks
for complexity to be \emph{determined}.
\end{conceptbox}

\section*{Coverage of the HEC 2025 outcomes}

HEC 2025 publishes no course contents for any course --- only outcomes. Its
five bullets for this course map onto the same chapters, so that a course file
written against either edition has something to check.

\rowcolors{2}{white}{lightbg}
\begin{longtable}{@{}L{7.6cm}L{2.6cm}L{4.8cm}@{}}
\toprule
\rowcolor{primary!15}
\textbf{HEC 2025 outcome} & \textbf{Chapters} & \textbf{Met by} \\
\midrule
\endhead
Design efficient algorithms for common computational problems & 7--9, 14--16,
18--21 & CLO-6 and CLO-7 \\
Analyze algorithm complexity using Big~O notation & 3, 5, 6 & CLO-5 \\
Solve problems involving recursion, divide-and-conquer and dynamic
programming & 5--9, 15, 16 & CLO-6 \\
Compare different algorithmic approaches to problem-solving & 7, 14--16, 22 &
CLO-6; the comparison tables that close Parts II and IV \\
Demonstrate correctness and optimality of algorithms & 4, 10, 14, 20 & Loop
invariants in 4; the lower bound in 10; exchange arguments in 14; Dijkstra's
correctness proof in 20 \\
\bottomrule
\end{longtable}

\begin{alertbox}[Where this course sits in the degree, and the one awkwardness]
HEC 2023 places this course in \textbf{Semester 7} and \emph{Artificial
Intelligence} in Semester~3 --- so a student meets heuristic search four
semesters before meeting the analysis that would let them reason about its
cost. HEC 2025 restores the sensible order, putting algorithms in Semester~IV
and AI in Semester~V.

\smallskip
If your department follows the 2023 scheme, this handout assumes nothing from
the later courses and can be taught as published. If it follows 2025, this
course comes earlier and Chapters 12 and 13 may need more time, because Data
Structures will be fresher but less practised.
\end{alertbox}

\newpage

%=====================================================================
\chapter*{The Course at a Glance}
\addcontentsline{toc}{chapter}{The Course at a Glance}
\markboth{The Course at a Glance}{}

\dsfigh{%
\begin{tikzpicture}[
  part/.style={rounded corners=3pt, fill=#1, text=white, font=\bfseries\footnotesize,
               minimum width=2.5cm, text width=2.2cm, minimum height=0.95cm,
               align=center},
  ch/.style={rounded corners=2pt, draw=#1, line width=0.7pt, fill=#1!7, text=#1!60!black,
             font=\scriptsize, minimum width=2.5cm, minimum height=0.52cm, align=center},
  arr/.style={-{Stealth[length=2mm]}, primary!45, line width=0.8pt}
]
\node[part=primary]   (p1) at (0,0)    {Part I\\Analysing};
\node[part=secondary] (p2) at (3.0,0)  {Part II\\Recursion and D\&C};
\node[part=greenhl]   (p3) at (6.0,0)  {Part III\\Sorting and Structures};
\node[part=purplehl]  (p4) at (9.0,0)  {Part IV\\Design Paradigms};
\node[part=tealhl]    (p5) at (12.0,0) {Part V\\Graphs};
\node[part=goldhl]    (p6) at (15.0,0) {Part VI\\Limits and Practice};

\foreach \a/\b in {p1/p2,p2/p3,p3/p4,p4/p5,p5/p6}{\draw[arr] (\a) -- (\b);}

\foreach \i/\t in {1/{1. Role and cost}, 2/{2. The three cases},
                   3/{3. Asymptotics}, 4/{4. Loop invariants}}{
  \pgfmathsetmacro{\yy}{-0.90-(\i-1)*0.62}
  \node[ch=primary] at (0,\yy) {\t};}

\foreach \i/\t in {1/{5. Recurrences I}, 2/{6. Master theorem},
                   3/{7. Brute force}, 4/{8. Merge sort}, 5/{9. Quicksort}}{
  \pgfmathsetmacro{\yy}{-0.90-(\i-1)*0.62}
  \node[ch=secondary] at (3.0,\yy) {\t};}

\foreach \i/\t in {1/{10. Lower bound}, 2/{11. Selection},
                   3/{12. Heaps, trees}, 4/{13. Hashing}}{
  \pgfmathsetmacro{\yy}{-0.90-(\i-1)*0.62}
  \node[ch=greenhl] at (6.0,\yy) {\t};}

\foreach \i/\t in {1/{14. Greedy}, 2/{15. DP: elements},
                   3/{16. DP: problems}, 4/{17. Amortised}}{
  \pgfmathsetmacro{\yy}{-0.90-(\i-1)*0.62}
  \node[ch=purplehl] at (9.0,\yy) {\t};}

\foreach \i/\t in {1/{18. Graphs}, 2/{19. MST},
                   3/{20. Shortest paths}}{
  \pgfmathsetmacro{\yy}{-0.90-(\i-1)*0.62}
  \node[ch=tealhl] at (12.0,\yy) {\t};}

\foreach \i/\t in {1/{21. Strings}, 2/{22. P and NP},
                   3/{23. Project}}{
  \pgfmathsetmacro{\yy}{-0.90-(\i-1)*0.62}
  \node[ch=goldhl] at (15.0,\yy) {\t};}

\node[rounded corners=3pt, fill=primary!7, draw=primary!30,
      minimum width=17.3cm, minimum height=0.66cm,
      font=\scriptsize\bfseries, text=primary] at (7.5,-4.35)
      {\faIcon{stopwatch}~~Every chapter: is it correct, what does it cost, and
       could anything do better?};
\end{tikzpicture}}%
{The six parts and the chapters in each. Chapter 10 is the hinge: it is the
first result in the book about a \emph{problem} rather than about an
algorithm.}{coursemap}

\vspace{4pt}
{\small\itshape\color{gray!70!black}
This course is the published continuation of \emph{Data Structures} and sits
beside \emph{Theory of Automata}. Appendix~F states what belongs to each.}

\newpage

%=====================================================================
\chapter*{Notation}
\addcontentsline{toc}{chapter}{Notation}
\markboth{Notation}{}

\rowcolors{2}{white}{lightbg}
\begin{longtable}{@{}L{3.0cm}L{12.2cm}@{}}
\toprule
\rowcolor{primary!15}
\textbf{Symbol} & \textbf{Meaning} \\
\midrule
\endhead
$n$ & The size of the input. Chapter 2 is about what that sentence hides \\
$T(n)$ & The running time on an input of size $n$, in the model of Chapter 1 \\
$\lg n$ & $\log_2 n$. Used throughout; $\ln$ is natural and $\log$ without a
base means the base does not matter \\
$\bigO{f(n)}$ & An asymptotic \emph{upper} bound: $T$ grows no faster than $f$ \\
$\bigOm{f(n)}$ & An asymptotic \emph{lower} bound \\
$\bigTh{f(n)}$ & Both: $T$ grows exactly as fast as $f$ \\
$o(f(n))$, $\omega(f(n))$ & Strict versions of $O$ and $\Omega$ --- strictly
slower, strictly faster \\
$A[1 \mathinner{\ldotp\ldotp} n]$ & An array of $n$ elements, indexed
\textbf{from 1}, as the pseudocode does. The C++ in \texttt{code/} indexes from
0, and says so where it matters \\
$|V|$, $|E|$ & The number of vertices and edges of a graph. Chapter 18 explains
why bounds are written in both \\
$\hat{c}_i$ & The amortised cost of the $i$th operation (Chapter 17) \\
$\Phi$ & A potential function (Chapter 17) \\
$\mathrm{P}$, $\mathrm{NP}$ & The complexity classes of Chapter 22 \\
$\le_p$ & Polynomial-time reducibility (Chapter 22) \\
$\lceil x \rceil$, $\lfloor x \rfloor$ & Ceiling and floor \\
\bottomrule
\end{longtable}

\begin{alertbox}[Three notational habits worth breaking now]
\textbf{$O$ is not $\Theta$.} ``Merge sort is $\bigO{n^2}$'' is \emph{true} and
nearly useless. If you mean the bound is tight, write $\Theta$. Chapter 3 is
blunt about this because the examination is.

\smallskip
\textbf{$=$ in ``$T(n) = \bigO{n^2}$'' is not equality.} It is set membership
by long-standing abuse: $O(n^2)$ is a \emph{set of functions} and $T$ belongs
to it. You may never write $\bigO{n^2} = T(n)$, and $\bigO{n} = \bigO{n^2}$ is
meaningless.

\smallskip
\textbf{Constants are discarded, not absent.} $\bigTh{n}$ with a constant of
$10^6$ loses to $\bigTh{n^2}$ with a constant of $1$ until $n$ exceeds a
million. The measured strand in every chapter exists to keep this in view.
\end{alertbox}

\newpage

%=====================================================================
\chapter*{Prerequisite Self-Check}
\addcontentsline{toc}{chapter}{Prerequisite Self-Check}
\markboth{Prerequisite Self-Check}{}

Twelve questions. If you can answer ten, begin at Chapter~1. If you cannot,
Appendix~B covers the mathematics and your Data Structures notes cover the
rest --- and the gap is worth closing in Week~1 rather than Week~9.

\begin{enumerate}[leftmargin=2.4em, itemsep=4pt, topsep=4pt]
  \item Write the sum $1 + 2 + \cdots + n$ in closed form, and say why.
        \hfill{\scriptsize\color{neutral}(Appendix~B)}
  \item Evaluate $\sum_{i=0}^{k} 2^{i}$. What is it approximately, for large
        $k$? \hfill{\scriptsize\color{neutral}(Appendix~B)}
  \item Simplify $\log_2(n^3)$ and $2^{\log_2 n}$.
        \hfill{\scriptsize\color{neutral}(Appendix~B)}
  \item State the principle of mathematical induction, and use it to show
        $\sum_{i=1}^{n} i = n(n+1)/2$.
        \hfill{\scriptsize\color{neutral}(Appendix~B)}
  \item A binary tree of height $h$ has at most how many leaves? Why?
        \hfill{\scriptsize\color{neutral}(Data Structures)}
  \item What is the height of a binary heap containing $n$ elements?
        \hfill{\scriptsize\color{neutral}(\chref{structures})}
  \item Give the worst-case height of an unbalanced binary search tree on $n$
        keys, and the input that produces it.
        \hfill{\scriptsize\color{neutral}(\chref{structures})}
  \item Describe how a hash table resolves collisions by chaining, and what
        happens as the load factor rises.
        \hfill{\scriptsize\color{neutral}(\chref{hashing})}
  \item What is the difference between an adjacency matrix and an adjacency
        list, in space? \hfill{\scriptsize\color{neutral}(\chref{graphs})}
  \item If a fair coin is tossed $n$ times, what is the expected number of
        heads? \hfill{\scriptsize\color{neutral}(Appendix~B)}
  \item In C++, what is the difference between \texttt{std::vector} and a raw
        array, and which one has $\bigO{1}$ indexing?
        \hfill{\scriptsize\color{neutral}(Appendix~A)}
  \item Write a function that reverses an array in place. What does it cost?
        \hfill{\scriptsize\color{neutral}(Programming Fundamentals)}
\end{enumerate}

\begin{conceptbox}
Notice what is \emph{not} on that list: calculus, linear algebra, and any
particular cleverness. This course needs fluency with sums, logarithms and
induction, comfort with the data structures you already built, and a
willingness to finish a proof rather than stop when the idea is clear.

\smallskip
It also needs you to be able to compile and run a C++ program. Appendix~A sets
that up in one sitting; do it before Week~2, because from Chapter~1 onward
every chapter asks you to measure something.
\end{conceptbox}

\newpage
\tableofcontents
\newpage
\listoffigures
\newpage

\pagenumbering{arabic}
